\subsection{Avaliação final dos experimentos}

Ao longo das sete aulas, foram produzidas e avaliadas amostras de sabão por saponificação a quente de uma mistura de óleo de coco babaçu e óleo de soja. As seis primeiras aulas foram dedicadas à produção, procurando-se variar um fator por vez: primeiro a proporção entre os óleos (aulas 1 e 2), depois a quantidade de ácido esteárico (aulas 3 e 4) e, por fim, a adição de etanol como cossolvente (aulas 5 e 6). Enquanto que a última aula foi destinada à caracterização dos produtos. Esta seção sintetiza os principais acertos, limitações e aprendizados desse percurso.

\subsubsection{O que deu certo}

A escolha da rota a quente mostrou-se acertada: o aquecimento reduziu o tempo de cura a cerca de uma semana, o que viabilizou a validação dos produtos das duas últimas aulas (5 e 6) na aula de caracterização sendo assim, mantendo as aulas dento do cronograma acadêmico — algo inviável pela rota a frio, que exigiria de quatro a seis semanas de cura.

A mistura de óleos também se mostrou uma boa estratégia, combinando os ácidos graxos saturados de cadeia média do babaçu, que conferem dureza e boa formação de espuma, com a maciez e o baixo custo do óleo de soja. A composição de aproximadamente 33\% de babaçu e 66\% de soja, obtida na amostra~1 do primeiro dia, gerou uma barra de boa qualidade e serviu de referência para as aulas seguintes. Manter a quantidade de água fixa em 25\% da massa de óleo ao longo de todas as práticas reduziu o número de variáveis livres e facilitou a comparação entre os experimentos.

A introdução do ácido esteárico, a partir da terceira aula, confirmou seu papel como agente estruturante: a barra atingiu, com apenas uma semana de cura, a firmeza que a amostra de referência só havia alcançado após duas semanas. Já o etanol, empregado nas aulas 5 e 6, foi o aditivo de grande impacto positivo também. Além de reduzir o pH de bancada (10 e 10,5, contra os 13 das primeiras aulas), melhorou visivelmente a aparência, a textura e até o odor do produto, tornando essas as melhores amostras do conjunto.

Por fim, a caracterização indicou que, após a cura completa, todas as amostras apresentaram pH em torno de 9 e formação de espuma abundante e estável ao longo de dez minutos e que se manteve até o teste de pH (durando mais de 35 minutos). Esse resultado é particularmente relevante, pois mostra que o pH residual elevado observado logo após a etapa de bancada foi efetivamente reduzido ao longo da cura, aproximando os produtos da faixa desejada. O teste de limpeza complementar, realizado em casa, confirmou ainda a capacidade dos sabões de remover gordura em água, com destaque para a amostra da quinta aula.

\subsubsection{O que deu errado}

Os principais problemas estiveram ligados ao controle da quantidade de base e à constância das condições. Na primeira aula, a substituição do óleo de coco pelo babaçu sem o recálculo da quantidade de NaOH resultou em excesso de base; na amostra~2, com maior teor de babaçu, esse excesso foi acentuado e levou a uma saponificação quase instantânea, gerando um produto granular e quebradiço, o pior resultado de todo o conjunto. Na segunda aula, o cálculo do NaOH foi feito para a mistura de babaçu e soja enquanto a proporção entre os óleos também era alterada, de modo que dois fatores variaram simultaneamente; somou-se a isso o uso de uma soda em lascas aparentemente contaminada por umidade, o que compromete a estimativa da base efetivamente reagindo. Por sorte a reação não se efetivou rápido o suficiente e o problema do uso reduzido hidróxido de sódio pode ser revertido antes do fim do experimento.

Outras falhas pontuais comprometeram a constância das condições: na quarta aula, o óleo foi superaquecido a mais de 110~°C por descuido, e na sexta aula, empregou-se acidentalmente ácido esteárico em quantidade abaixo da planejada, o que só foi percebido durante a escrita do relatório, e a soda cáustica, já no fim do estoque, produziu uma solução com impurezas visíveis.

A própria etapa de caracterização apresentou limitações. A ausência de um sabão-controle dificultou uma avaliação mais palpável dos produtos. Além disso, o método de determinação da densidade mostrou-se impreciso: o deslocamento medido foi de apenas 2,0 a 2,5~ml em provetas com baixa resolução, de modo que pequenas variações de leitura alteram bastante o valor calculado (Tabela~\ref{tab:dia7_caracterizacao}). Esse problema fica evidente em uma incoerência observada, os sabões das duas primeiras aulas, apesar de apresentarem densidade calculada inferior à da água, afundaram quando imersos, o que sugere um erro de medida nessas duas amostras. Por fim, a amostra da segunda aula apresentou a maior solubilidade, o que pode estar associado a uma saponificação incompleta — seu interior continha uma quantidade de bolhas de óleo não saponificado bem maior que a das demais.

\subsubsection{Melhor e pior produto}

Considerando aparência, textura, odor, pH e, sobretudo, poder de limpeza, o melhor produto foi o da quinta aula, primeira a empregar etanol. Além de combinar boa consistência, coloração agradável e o menor pH de bancada, foi também o que melhor dissolveu a gordura no teste de limpeza (Figura~\ref{fig:limpeza_resultados}). Em contraposição, o pior foi a amostra~2 da primeira aula, cuja saponificação descontrolada, decorrente do grande excesso de base, resultou em um material granular.
